% !TeX program = xelatex
% ============================================================
% project_description/project_description.tex — «Детальное описание
% проекта» к КТ1/КТ2 курсового проекта ОП «Программная инженерия».
%
% Обязательный деливерабл КТ1 (обновляется на КТ2) по Программе практики:
%   • прикладной вид  — Приложение 6 (планы/этапы, тех.стек+обоснование,
%     >=3 количественных критериев оценки, особые пометки);
%   • исследовательский вид — Приложение 7 (цель, мотивация, объект,
%     предмет, гипотеза, обзор, планы/этапы, >=3 критериев, пометки).
% Структура ветвится по project.kind, язык — по project.lang.
% Заполняйте жёлтые \hseFill{…}; данные о себе и проекте — в настройках.
% ============================================================
\documentclass[12pt,a4paper]{extarticle}
\input{preamble}
\begin{document}
\sloppy

% ── Шапка документа ──
\begin{center}

{\large\bfseries Detailed Project Description}\\[2pt]
{\normalsize (Checkpoints 1--2, course project)}\\[6pt]

\textbf{\projectname}\\[4pt]

\hseAuthorName, group \hseGroup

\end{center}
\vspace{0.4cm}

((* if project.kind == 'research' *))
% ============================================================
%  ПРИЛОЖЕНИЕ 7 — исследовательский проект
% ============================================================

\section{Main plans and stages of the project}
\subsection{Project description}
\textbf{Aim of the project:} \hseFill{one sentence — the scientific result to be achieved}.

\textbf{Motivation} (why this project): \hseFill{why the topic is relevant and interesting}.

\textbf{Object of study:} \hseFill{object}.

\textbf{Subject of study:} \hseFill{subject}.

\textbf{Hypothesis (brief):} \hseFill{the hypothesis to be tested}.

\textbf{Domain overview:} \hseFill{a short overview of the domain and existing approaches}.

\subsection{Plans and stages}
\begin{longtable}{|>{\raggedright\arraybackslash}p{3cm}|>{\raggedright\arraybackslash}p{5cm}|>{\raggedright\arraybackslash}p{3.5cm}|>{\raggedright\arraybackslash}p{2cm}|}
\hline
\textbf{Stage} & \textbf{Description of work} & \textbf{Expected results} & \textbf{Due date} \\ \hline
\endhead
\hseFill{stage 1} & \hseFill{what is done} & \hseFill{result} & \hseFill{date} \\ \hline
\hseFill{stage 2} & \hseFill{what is done} & \hseFill{result} & \hseFill{date} \\ \hline
\end{longtable}

\section{Evaluation criteria}
Provide at least three quantitative criteria by which the project result will be assessed (e.g. novelty of the study, quality of the literature review, correctness of the model/algorithms, quality of the implementation, depth of the analysis of experimental results).
\begin{longtable}{|>{\raggedright\arraybackslash}p{5cm}|>{\raggedright\arraybackslash}p{11cm}|}
\hline
\textbf{Criterion} & \textbf{Description} \\ \hline
\endhead
\hseFill{criterion 1} & \hseFill{how it is measured} \\ \hline
\hseFill{criterion 2} & \hseFill{how it is measured} \\ \hline
\hseFill{criterion 3} & \hseFill{how it is measured} \\ \hline
\end{longtable}

\section{Special notes}
\hseFill{any risks, changes to the plans, and other important details for this stage}.

((* else *))
% ============================================================
%  ПРИЛОЖЕНИЕ 6 — прикладной (программный) проект
% ============================================================

\section{Main plans and stages of the project}
\subsection{Brief description of the project}
\textbf{Project name:} \projectname.

\textbf{Aim of the project:} \hseFill{one sentence — what is built}.

\textbf{Brief description of the tasks:} \hseFill{the key tasks of the project}.

\subsection{Plans and stages}
\begin{longtable}{|>{\raggedright\arraybackslash}p{3cm}|>{\raggedright\arraybackslash}p{5cm}|>{\raggedright\arraybackslash}p{3.5cm}|>{\raggedright\arraybackslash}p{2cm}|}
\hline
\textbf{Stage} & \textbf{Description of work} & \textbf{Expected results} & \textbf{Due date} \\ \hline
\endhead
\hseFill{stage 1} & \hseFill{what is done} & \hseFill{result} & \hseFill{date} \\ \hline
\hseFill{stage 2} & \hseFill{what is done} & \hseFill{result} & \hseFill{date} \\ \hline
\end{longtable}

\section{Technology stack and its justification}
\subsection{List of technologies used}
\begin{longtable}{|>{\raggedright\arraybackslash}p{4cm}|>{\raggedright\arraybackslash}p{5cm}|>{\raggedright\arraybackslash}p{7cm}|}
\hline
\textbf{Technology / tool} & \textbf{Description} & \textbf{Reasons for the choice} \\ \hline
\endhead
\hseFill{technology 1} & \hseFill{what it is} & \hseFill{why chosen} \\ \hline
\hseFill{technology 2} & \hseFill{what it is} & \hseFill{why chosen} \\ \hline
\end{longtable}

\subsection{Justification of the chosen stack}
\hseFill{why this stack was chosen, what advantages it has for this project, and how it helps achieve the project goals}.

\section{Evaluation criteria}
Provide at least three quantitative criteria by which the project result will be assessed (e.g. \% of functional requirements met, number of implemented functions, response time, test coverage \%, number of bugs per KLOC, cyclomatic complexity, number of supported platforms).
\begin{longtable}{|>{\raggedright\arraybackslash}p{5cm}|>{\raggedright\arraybackslash}p{11cm}|}
\hline
\textbf{Criterion} & \textbf{Description} \\ \hline
\endhead
\hseFill{criterion 1} & \hseFill{how it is measured} \\ \hline
\hseFill{criterion 2} & \hseFill{how it is measured} \\ \hline
\hseFill{criterion 3} & \hseFill{how it is measured} \\ \hline
\end{longtable}

\section{Special notes}
\hseFill{any risks, changes to the plans, and other important details for this stage}.

((* endif *))

\end{document}
